\section{Conclusions}
In summary, the T2K experiment has measured charged current inclusive events, in a restricted phase
space of $\theta_{\mu}<32^{\circ}$ and $p_{\mu}>$500 MeV/c,
the flux averaged cross sections (cm$^{2}$ per nucleon) and ratio
of cross sections, as
\begin{equation}
\sigma(\overline{\nu})=\left(0.900\pm0.029{\rm (stat.)}\pm0.088{\rm (syst.)}\right)\times10^{-39},\label{eq:conclusion-1}
\end{equation}
\begin{equation}
\sigma(\nu)=\left(2.41\pm0.021{\rm (stat.)}\pm0.231{\rm (syst.)}\ \right)\times10^{-39}\label{eq:conclusion-2}
\end{equation}
and
\begin{equation}
R\left(\frac{\sigma(\overline{\nu})}{\sigma(\nu)}\right)=0.373\pm0.012{\rm (stat.)}\pm0.015{\rm (syst.).}\label{eq:conclusion-3}
\end{equation}
The \nubar inclusive cross section and the ratio $R$ results
are the first published measurements at \nuandnubar flux energies\citep{flux-model} below
1.5 GeV.  Although the current uncertainty on the different model combinations is relatively large, 
we expect future higher statistics comparisons will be valuable for model discrimination.
